\documentclass[aspectratio=169,10pt,english]{beamer}

\input{../../../_preambulo_beamer.tex}
\input{../../../_comandos_beamer.tex}

\selectlanguage{english}

\title{MA1001B - Statistical Modeling for Decision Making}
\subtitle{Lesson 35: t Test for a Mean with Unknown Variance}
\author{}
\institute{Tecnol\'ogico de Monterrey}
\date{}

\begin{document}

\begin{frame}[plain]
  \vspace{1.2cm}
  {\Large\bfseries MA1001B - Statistical Modeling for Decision Making\par}
  \vspace{0.7cm}
  {\large Lesson 35: t Test for a Mean with Unknown Variance\par}
  \vspace{0.7cm}
  {\color{TecRojo}\rule{\textwidth}{1.2pt}}
  \vspace{0.8cm}

  MA1001B Faculty\\[0.3cm]
  Term\\[0.3cm]
  Tecnol\'ogico de Monterrey
\end{frame}

\begin{frame}{The Unknown-Sigma Mean Question}
  \begin{alertblock}{Guiding question}
    How do we test $H_0:\mu=50$ when only the sample standard deviation $s$ is
    available?
  \end{alertblock}

  \vspace{0.6em}
  By the end of this lesson, you can:
  \begin{itemize}
    \item replace a known-sigma $z$ with Student's $t$ and $df=n-1$;
    \item compute $t=(\bar{x}-\mu_0)/(s/\sqrt{n})$;
    \item read a two-sided p-value and decide at $\alpha=0.05$;
    \item explain why $n=8$ plus skew makes that p-value fragile.
  \end{itemize}

  \vfill
  \scriptsize All pack-out cycle times used today are synthetic. No real company data are used.
\end{frame}

\begin{frame}{From $z$ to $t$ when $\sigma$ Is Unknown}
  \small
  A synthetic pack-out line claims a mean cycle time of $50$ minutes.
  Operations samples $n=40$ completed packs. The population $\sigma$ is not
  given, so the standard error uses $s$ and the reference curve is $t_{n-1}$.

  \[
    t=\frac{\bar{x}-\mu_0}{s/\sqrt{n}},\qquad df=n-1.
  \]

  \begin{exampleblock}{Two-sided hypotheses}
    $H_0:\mu=50$ versus $H_1:\mu\neq 50$, with $\alpha=0.05$. Both tails of
    $t_{39}$ count toward the p-value.
  \end{exampleblock}
\end{frame}

\begin{frame}[fragile]{Step 1: Sample $s$ and the $t$ Statistic}
  \begin{columns}[T]
    \begin{column}{0.42\textwidth}
      \small
      \begin{lstlisting}
se = s / sqrt(n)
t = (xbar - 50) / se
df = n - 1
      \end{lstlisting}

      \begin{block}{Verified output}
        $\bar{x}=52.229590$\\
        $s=4.955099$, $SE=0.783470$\\
        $df=39$, $t=2.845789$
      \end{block}
    \end{column}
    \begin{column}{0.58\textwidth}
      \centering
      \includegraphics[width=\textwidth]{../figures/t_test_mean_01_sample_standard_error.png}
      \vspace{0.2em}
      \scriptsize Histogram of $n=40$ pack times from Step 1
    \end{column}
  \end{columns}
\end{frame}

\begin{frame}{Two-Sided p-Value from $t_{39}$}
  \small
  The observed statistic is compared with Student $t$ on $39$ degrees of
  freedom, not with a standard normal:
  \[
    p=2P(T_{39}\ge 2.845789)=0.007026.
  \]
  The two-sided critical value at $\alpha=0.05$ is
  \[
    t^{*}=t_{39,0.975}=2.022691.
  \]
  Because $|t|>t^{*}$ and $p<\alpha$, the test rejects $H_0:\mu=50$.
\end{frame}

\begin{frame}[fragile]{Step 2: Two-Sided Decision}
  \begin{columns}[T]
    \begin{column}{0.40\textwidth}
      \small
      \begin{lstlisting}
p = 2 * t.sf(abs(t), df)
t_star = t.ppf(0.975, df)
      \end{lstlisting}

      \begin{block}{Verified output}
        $p=0.007026$ (manual and scipy)\\
        $t^{*}=2.022691$\\
        Decision: reject $H_0$
      \end{block}
    \end{column}
    \begin{column}{0.60\textwidth}
      \centering
      \includegraphics[width=\textwidth]{../figures/t_test_mean_02_two_sided_pvalue.png}
      \vspace{0.2em}
      \scriptsize Two-sided $t$ tails from Step 2
    \end{column}
  \end{columns}
\end{frame}

\begin{frame}{Step 3: $n=8$ plus Skew Makes $p$ Fragile}
  \begin{columns}[T]
    \begin{column}{0.42\textwidth}
      \small
      A rushed $n=8$ follow-up has skewness $1.935033$ and Shapiro-Wilk
      $p=0.001286$.

      \vspace{0.4em}
      Full sample: $t=0.307397$, $p=0.767485$.

      \vspace{0.4em}
      Drop the $83.174379$-minute delay: $p=0.081584$.

      \vspace{0.5em}
      \textbf{Limit:} with small $n$ and a long tail, one observation moves
      $p$ by an order of magnitude. The $n=40$ p-value stays between
      $0.001481$ and $0.013087$.
    \end{column}
    \begin{column}{0.58\textwidth}
      \centering
      \includegraphics[width=\textwidth]{../figures/t_test_mean_03_small_n_skew_limit.png}
      \vspace{0.2em}
      \scriptsize Leave-one-out p-values from Step 3
    \end{column}
  \end{columns}
\end{frame}

\begin{frame}{Concept Check}
  \begin{enumerate}
    \item Why is the reference curve $t_{39}$ rather than $N(0,1)$?
    \item Compute $SE=s/\sqrt{n}$ from $s=4.955099$ and $n=40$.
    \item If $p=0.007026$ and $\alpha=0.05$, what is the two-sided decision?
    \item Why does dropping one long delay change the $n=8$ p-value so much?
  \end{enumerate}

  \vspace{0.6em}
  \begin{block}{Connection to Lesson 36}
    The session can continue directly into a hypothesis test for a business
    proportion, where the standard error uses $p_0$ rather than $\hat{p}$.
  \end{block}

  \vfill
  \scriptsize Source: OpenStax, \textit{Introductory Business Statistics 2e},
  Chapter 9, Sections 9.3--9.4. CC BY-NC-SA.
\end{frame}

\appendix



\end{document}
